% !TeX root = ../../Book2.tex
\section{极值情形}
\label{b2:sec:2002}

Heisenberg 不等式的证明经过两个估计：先取复内积的实部，再用向量 Cauchy--Schwarz。要识别极值函数，必须同时追踪这两步的等号条件。这样得到的微分方程还说明了两种不同的损失：不同方向的宽度不一致会破坏向量比例关系，二次相位则会使位置与频率产生相关。

\subsection{等号情形}
\label{b2:subsec:200201}

\Needspace{6\baselineskip}
\begin{proposition}[Heisenberg 等号情形]\label{b2:thm:20-equality}
设 \(d\geq1\)，\(f\in L^2(\mathbb R^d;\mathbb C)\)、\(\|f\|_2=1\)，且 \(|x|f\)、\(|\xi|\widehat f\) 都属于 \(L^2\)，其中 \(\widehat f\) 为负号 \(2\pi\) 归一化的 Plancherel 变换。定义
\[
 \begin{aligned}
 a&=\int_{\mathbb R^d}x|f(x)|^2\,\mathrm dx,
 \quad
 b=\int_{\mathbb R^d}\xi|\widehat f(\xi)|^2\,\mathrm d\xi,\\
 \sigma_x^2&=\int_{\mathbb R^d}|x-a|^2\cdot|f(x)|^2\,\mathrm dx,\\
 \sigma_\xi^2&=\int_{\mathbb R^d}|\xi-b|^2\cdot
       |\widehat f(\xi)|^2\,\mathrm d\xi.
 \end{aligned}
\]
则 \(\sigma_x\sigma_\xi=d/(4\pi)\)，当且仅当存在 \(a_0>0\) 和常数 \(C\in\mathbb C\)，使
\[
 f(x)=C e^{-\pi a_0|x-a|^2}e^{2\pi\mathrm i b\cdot x}
 \quad
 |C|=(2a_0)^{d/4}
\]
几乎处处成立。
\end{proposition}
\begin{proof}
先把两个中心移到零。令
\[
 g(y)=e^{-2\pi\mathrm i b\cdot(y+a)}\cdot f(y+a).
\]
平移与调制保持 \(L^2\) 范数，并将位置中心和频率中心分别移到零。有限频率二阶矩保证 \(g\in H^1(\mathbb R^d)\)，且
\[
 \||y|g\|_2=\sigma_x,
 \quad
 \|\nabla g\|_2=2\pi\sigma_\xi.
\]
在复 Hilbert 空间 \(L^2(\mathbb R^d;\mathbb C^d)\) 中，取
\[
 X(y)=y\cdot g(y),
 \quad
 Y(y)=\nabla g(y),
 \quad
 \langle X,Y\rangle
 =\int_{\mathbb R^d}X(y)\cdot\overline{Y(y)}\,\mathrm dy.
\]
这里内积对第一个变量线性。由\cref{b2:lem:20-weighted-smooth-density}，位置权重与弱梯度的平方可积性允许积分分部，得到
\[
 \operatorname{Re}\langle X,Y\rangle=-\frac d2.
\]
因此 \(X,Y\) 均非零，原不等式恰来自
\[
 \frac d2
 =-\operatorname{Re}\langle X,Y\rangle
 \leq|\langle X,Y\rangle|
 \leq\|X\|_2\cdot\|Y\|_2.
\]

若取等，第二步的等号给出 \(Y=\lambda X\)，其中 \(\lambda\) 是一个常复数。这是整个向量 Hilbert 空间中的线性相关，不是允许每个位置或每个坐标分别选择比例因子。确切地说，令
\[
 \lambda=\frac{\langle Y,X\rangle}{\|X\|_2^2},
\]
则
\[
 \|Y-\lambda X\|_2^2
 =\|Y\|_2^2-\frac{|\langle X,Y\rangle|^2}{\|X\|_2^2}=0.
\]
第一步的等号又要求 \(\langle X,Y\rangle\) 为负实数。由于
\(\langle X,Y\rangle=\overline\lambda\|X\|_2^2=-d/2\)，存在 \(c>0\)，使
\[
 \nabla g(y)=-c\,y\cdot g(y)
\]
在弱导数意义下成立。

这个弱方程已经足以确定函数，无须先假设 \(g\) 经典可微。令
\[
 u(y)=e^{c|y|^2/2}\cdot g(y).
\]
在每个有界球内，指数及其一阶导数有界，故 \(u\in H^1_{\mathrm{loc}}\)，弱乘积公式给出 \(\nabla u=0\)。固定一个球，在稍大的球内将 \(u\) 磨光；所得光滑函数在原球内梯度为零，因此为常数。磨光函数在 \(L^1\) 局部收敛到 \(u\)，这些常数也收敛，因为两个常数之差的 \(L^1\) 范数等于球的体积乘其绝对值。于是 \(u\) 在每个球上几乎处处为常数。取递增的同心球，常数在交叠处一致，得到
\[
 g(y)=C e^{-c|y|^2/2}.
\]
令 \(a_0=c/(2\pi)\)，再恢复平移与调制，就是所述形式。由
\[
 \int_{\mathbb R^d}e^{-2\pi a_0|y|^2}\,\mathrm dy
 =(2a_0)^{-d/2}
\]
和 \(\|g\|_2=1\)，可得 \(|C|=(2a_0)^{d/4}\)。

反过来，所述 Gaussian 具有有限的两个二阶矩。中心化后有
\(\nabla g=-2\pi a_0\,y g\)，故上面两步估计都取等；它的位置中心为零，频率密度也是关于零对称的 Gaussian。恢复平移与调制后，两中心正是 \(a,b\)，方差乘积仍为 \(d/(4\pi)\)。
\end{proof}

极值族允许任意位置中心、频率中心、整体相位和共同宽度，却不允许不同方向使用不同的宽度。后一个限制来自同一个向量比例常数。下一组计算会把这种限制写成协方差矩阵的关系。

\subsection{Gaussian 函数}
\label{b2:subsec:200202}

先看每个方向宽度相同的情形。对 \(t>0\)，令
\[
 f_t(x)=(2t)^{d/4}e^{-\pi t|x|^2}.
\]
一维 Gaussian 积分的乘积给出 \(\|f_t\|_2=1\)，而 Gaussian 变换公式给出
\[
 \widehat f_t(\xi)
 =(2t)^{d/4}t^{-d/2}e^{-\pi|\xi|^2/t}
 =f_{1/t}(\xi).
\]
两个密度均为偶函数，故中心为零。为计算方差，对
\[
 \int_{\mathbb R}e^{-\pi s u^2}\,\mathrm du=s^{-1/2}
\]
关于 \(s>0\) 求导，得到
\[
 \int_{\mathbb R}u^2e^{-\pi s u^2}\,\mathrm du
 =\frac1{2\pi s^{3/2}}.
\]
求导的合法性可在任意固定 \(s_0>0\) 的小邻域核验：被积函数导数的绝对值由常数倍
\(u^2e^{-\pi s_0u^2/2}\) 控制，后者可积。代入 \(s=2t\)，再对其余坐标积分，每个位置坐标的方差都是 \(1/(4\pi t)\)。利用 \(\widehat f_t=f_{1/t}\)，得到
\[
 \sigma_x^2=\frac d{4\pi t},
 \quad
 \sigma_\xi^2=\frac{dt}{4\pi},
 \quad
 \sigma_x\sigma_\xi=\frac d{4\pi}.
\]
增大 \(t\) 会缩窄位置密度，同时加宽频率密度；二者的乘积保持不变。

不同方向的宽度由实对称正定矩阵描述。若非负密度 \(p\) 满足 \(\int p=1\)、\(\int xp=0\) 及 \(\int|x|^2p(x)\,\mathrm dx<\infty\)，它的\term[xiefangcha-juzhen]{协方差矩阵}[covariance matrix]是
\(\int xx^{\mathsf T}p(x)\,\mathrm dx\)：对角元记录各坐标的方差，非对角元记录位置坐标之间的相关。这里的矩阵与稍后涉及位置和频率的相关量用途不同。

\Needspace{6\baselineskip}
\begin{proposition}[Gaussian 的协方差]\label{b2:prop:20-gaussian-covariance}
设 \(d\geq1\)，\(A\) 是 \(d\times d\) 实对称正定矩阵，令
\[
 f_A(x)=\det(2A)^{1/4}e^{-\pi x^{\mathsf T}Ax},
 \quad
 x\in\mathbb R^d.
\]
则 \(f_A\) 的 \(L^2\) 范数为一，其负号 \(2\pi\) 归一化 Fourier 变换为
\[
 \widehat f_A(\xi)
 =\det(2A)^{1/4}(\det A)^{-1/2}
       e^{-\pi\xi^{\mathsf T}A^{-1}\xi}.
\]
位置中心和频率中心均为零，协方差矩阵满足
\[
 \begin{aligned}
 \Sigma_x
 &\coloneqq\int_{\mathbb R^d}xx^{\mathsf T}|f_A(x)|^2\,\mathrm dx
 =\frac1{4\pi}A^{-1},\\
 \Sigma_\xi
 &\coloneqq\int_{\mathbb R^d}\xi\xi^{\mathsf T}
                |\widehat f_A(\xi)|^2\,\mathrm d\xi
 =\frac1{4\pi}A.
 \end{aligned}
\]
因而总方差满足
\[
 \sigma_x^2\sigma_\xi^2
 =\frac{\operatorname{tr}A\cdot\operatorname{tr}A^{-1}}{16\pi^2}
 \geq\frac{d^2}{16\pi^2},
\]
等号当且仅当 \(A=tI\)，其中 \(t>0\)。
\end{proposition}
\begin{proof}
实对称矩阵的正交对角化给出
\[
 A=O^{\mathsf T}\Lambda O,
 \quad
 \Lambda=\operatorname{diag}(\lambda_1,\ldots,\lambda_d),
 \quad
 \lambda_j>0,
 \quad
 O^{\mathsf T}O=I.
\]
在变量 \(y=Ox\) 下，Lebesgue 测度不变，二次型成为
\(\sum_j\lambda_jy_j^2\)。故
\[
 \int_{\mathbb R^d}e^{-2\pi x^{\mathsf T}Ax}\,\mathrm dx
 =\prod_{j=1}^d(2\lambda_j)^{-1/2}
 =\det(2A)^{-1/2},
\]
这正好解释归一化因子的选择。

Fourier 积分绝对收敛。在同一换元下，\(x\cdot\xi=y\cdot(O\xi)\)，所以逐坐标应用一维 Gaussian 变换公式，得到
\[
 \begin{aligned}
 \int_{\mathbb R^d}
 e^{-\pi x^{\mathsf T}Ax}e^{-2\pi\mathrm i x\cdot\xi}\,\mathrm dx
 &=\prod_{j=1}^d
   \left[\lambda_j^{-1/2}
   e^{-\pi(O\xi)_j^2/\lambda_j}\right]\\
 &=(\det A)^{-1/2}e^{-\pi\xi^{\mathsf T}A^{-1}\xi}.
 \end{aligned}
\]
乘回振幅即得变换公式；其振幅也等于
\(\det(2A^{-1})^{1/4}\)，因此 \(\widehat f_A=f_{A^{-1}}\)。

在 \(y\) 坐标中，位置密度是
\[
 \prod_{j=1}^d
 \left[(2\lambda_j)^{1/2}e^{-2\pi\lambda_jy_j^2}\right].
\]
每个因子积分为一且为偶函数，因此一阶矩及不同坐标之间的混合二阶矩为零。前面的矩计算给出第 \(j\) 个对角元为 \(1/(4\pi\lambda_j)\)。换回 \(x=O^{\mathsf T}y\)，得到
\[
 \Sigma_x
 =\frac1{4\pi}O^{\mathsf T}\Lambda^{-1}O
 =\frac1{4\pi}A^{-1}.
\]
将同一公式用于 \(\widehat f_A=f_{A^{-1}}\)，便得到
\(\Sigma_\xi=(4\pi)^{-1}A\)。

总方差等于协方差矩阵的迹，故只需估计
\[
 \left(\sum_{j=1}^d\lambda_j\right)
 \cdot
 \left(\sum_{j=1}^d\lambda_j^{-1}\right)
 \geq
 \left(\sum_{j=1}^d
       \sqrt{\lambda_j}\cdot\lambda_j^{-1/2}\right)^2
 =d^2.
\]
这是有限维 Cauchy--Schwarz；等号要求两个正向量
\((\sqrt{\lambda_j})_j\) 与 \((\lambda_j^{-1/2})_j\) 成比例，即所有 \(\lambda_j\) 相同。正交对角化随即给出 \(A=tI\)。
\end{proof}

例如在二维取 \(A=\operatorname{diag}(t,t^{-1})\)，则
\[
 \Sigma_x=\frac1{4\pi}
 \begin{pmatrix}t^{-1}&0\\0&t\end{pmatrix},
 \quad
 \Sigma_\xi=\frac1{4\pi}
 \begin{pmatrix}t&0\\0&t^{-1}\end{pmatrix},
 \quad
 \sigma_x\sigma_\xi=\frac{t+t^{-1}}{4\pi}.
\]
每一个坐标方向都达到自己的一维下界，但全维下界 \(1/(2\pi)\) 只有在 \(t=1\) 时达到。原因是求总方差时先分别相加，再取乘积：一个方向在位置上窄、另一个方向在频率上窄，并不能让两个总方差同时最小。

再看二次相位。它不改变位置密度，却会改变弱导数，因而改变频率方差。一维中可以保留复配对的虚部，得到精确记录这一增量的不等式。

\Needspace{6\baselineskip}
\begin{proposition}[含相关项的不确定性不等式]\label{b2:prop:20-correlated-uncertainty}
设 \(f\in L^2(\mathbb R;\mathbb C)\)、\(\|f\|_2=1\)，且 \(xf\)、\(\xi\widehat f\in L^2\)，其中 \(\widehat f\) 为负号 \(2\pi\) 归一化的 Plancherel 变换。令
\[
 \begin{aligned}
 a&=\int_{\mathbb R}x|f(x)|^2\,\mathrm dx,
 \quad
 b=\int_{\mathbb R}\xi|\widehat f(\xi)|^2\,\mathrm d\xi,\\
 \sigma_x^2&=\int_{\mathbb R}(x-a)^2|f(x)|^2\,\mathrm dx,\\
 \sigma_\xi^2&=\int_{\mathbb R}(\xi-b)^2
                    |\widehat f(\xi)|^2\,\mathrm d\xi.
 \end{aligned}
\]
有限频率二阶矩使弱导数 \(f'\) 属于 \(L^2\)。定义
\[
 Df=\frac1{2\pi\mathrm i}f',
 \quad
 \kappa=\operatorname{Re}\int_{\mathbb R}
       (x-a)\cdot f(x)\,\overline{(D-b)f(x)}\,\mathrm dx.
\]
则
\[
 \sigma_x^2\sigma_\xi^2
 \geq\kappa^2+\frac1{16\pi^2}.
\]
等号当且仅当存在 \(a_0>0\)、\(c\in\mathbb R\) 及常数 \(C\in\mathbb C\)，使
\[
 f(x)=C e^{-\pi(a_0+\mathrm ic)(x-a)^2}
           e^{2\pi\mathrm i bx},
 \quad
 |C|=(2a_0)^{1/4}
\]
几乎处处成立。对这些函数，
\[
 \sigma_x^2=\frac1{4\pi a_0},
 \quad
 \sigma_\xi^2=\frac{a_0^2+c^2}{4\pi a_0},
 \quad
 \kappa=-\frac{c}{4\pi a_0}.
\]
\end{proposition}
\begin{proof}
令 \(g(y)=e^{-2\pi\mathrm i b(y+a)}\cdot f(y+a)\)。它的两个中心为零，且
\[
 \|yg\|_2=\sigma_x,
 \quad
 \|Dg\|_2=\sigma_\xi,
 \quad
 \operatorname{Re}\int_{\mathbb R}yg\,\overline{Dg}\,\mathrm dy
 =\kappa.
\]
第二个等式使用 Plancherel 和 \(\widehat{Dg}(\xi)=\xi\widehat g(\xi)\)。所有配对都由 Cauchy--Schwarz 绝对可积。设
\[
 I=\int_{\mathbb R}y\cdot g(y)\cdot\overline{g'(y)}\,\mathrm dy.
\]
对 \(g\) 应用\cref{b2:lem:20-weighted-smooth-density}的一维情形，得到 \(\operatorname{Re}I=-1/2\)。注意复共轭会改变 \(D\) 中虚数因子的符号，因此
\[
 J\coloneqq\int_{\mathbb R}yg\,\overline{Dg}\,\mathrm dy
 =\frac{\mathrm i}{2\pi}I
 =\kappa-\frac{\mathrm i}{4\pi}.
\]
对 \(yg\) 和 \(Dg\) 应用复 Cauchy--Schwarz，得到
\[
 \kappa^2+\frac1{16\pi^2}
 =|J|^2
 \leq\|yg\|_2^2\cdot\|Dg\|_2^2
 =\sigma_x^2\sigma_\xi^2.
\]

由于 \(\operatorname{Im}J\neq0\)，两个向量均非零，等号当且仅当存在常复数 \(z\)，使 \(Dg=z\,yg\)。此时
\[
 \overline z\,\sigma_x^2=J
 =\kappa-\frac{\mathrm i}{4\pi},
 \quad
 \operatorname{Re}z=\frac{\kappa}{\sigma_x^2},
 \quad
 \operatorname{Im}z=\frac1{4\pi\sigma_x^2}>0.
\]
将 \(z=-c+\mathrm i a_0\) 写开，其中 \(a_0>0\)、\(c\in\mathbb R\)，弱微分方程变成
\[
 g'(y)=2\pi\mathrm i z\,y\cdot g(y).
\]
在每个有界区间中乘以光滑函数 \(e^{-\pi\mathrm i zy^2}\)，弱乘积公式说明 \(e^{-\pi\mathrm i zy^2}\cdot g(y)\) 的弱导数为零。用\cref{b2:thm:20-equality}证明中的局部磨光论证识别常数，得到
\[
 g(y)=C e^{\pi\mathrm i zy^2}
     =C e^{-\pi(a_0+\mathrm ic)y^2}.
\]
单位范数给出 \(|C|=(2a_0)^{1/4}\)，恢复中心后即得所述形式。

反过来，对这个归一化 Gaussian，
\[
 Dg=(-c+\mathrm i a_0)\,yg,
 \quad
 |g(y)|^2=(2a_0)^{1/2}e^{-2\pi a_0y^2}.
\]
其位置中心为零；由 Plancherel，
\[
 \int_{\mathbb R}\xi|\widehat g(\xi)|^2\,\mathrm d\xi
 =\int_{\mathbb R}Dg(y)\cdot\overline{g(y)}\,\mathrm dy
 =(-c+\mathrm i a_0)\int_{\mathbb R}y|g(y)|^2\,\mathrm dy=0.
\]
位置矩计算给出 \(\sigma_x^2=1/(4\pi a_0)\)。再取 \(Dg\) 的平方范数及它与 \(yg\) 的配对，得到
\[
 \sigma_\xi^2
 =|{-c+\mathrm i a_0}|^2\cdot\sigma_x^2
 =\frac{a_0^2+c^2}{4\pi a_0},
 \quad
 \kappa=-c\,\sigma_x^2=-\frac{c}{4\pi a_0}.
\]
这三个量满足所述等式。平移与调制恢复 \(a,b\)，不改变这些中心化量，证明完成。
\end{proof}

因此带二次相位的 Gaussian 都使加强式取等，而原来的 Heisenberg 下界还要求 \(\kappa=0\)，即 \(c=0\)。这里的相关量由两个 \(L^2\) 向量的复配对定义；位置密度和频率密度各在自己的变量空间中，不能把它们相乘来代替这个配对。

\subsection{Fourier 自对偶性}
\label{b2:subsec:200203}

本卷归一化下，\(e^{-\pi|x|^2}\) 的 Fourier 变换仍为自身。上面的归一化族满足 \(\widehat f_t=f_{1/t}\)，所以 \(t=1\) 恰是这个族的自对偶宽度；但每个 \(t>0\) 都是 Heisenberg 极值。自对偶要求变换前后函数相同，极值只要求位置与频率宽度的乘积最小，这两个条件并不等价。

反方向的区别也可用一个显式函数看到。只在一维，记 \(q(x)=e^{-\pi x^2}\)。它属于 Schwartz 空间且 \(\widehat q=q\)，故可逐次微分 Fourier 积分，得到
\[
 \widehat{x^nq}(\xi)
 =\left(\frac{\mathrm i}{2\pi}\right)^n q^{(n)}(\xi).
\]
对 \(q\) 求二阶和四阶导数，可直接算得
\[
 \begin{aligned}
 \widehat{x^2q}(\xi)
 &=\left(\frac1{2\pi}-\xi^2\right)q(\xi),\\
 \widehat{x^4q}(\xi)
 &=\left(\xi^4-\frac3\pi\xi^2+\frac3{4\pi^2}\right)q(\xi).
 \end{aligned}
\]
因此对
\[
 H(x)=(16\pi^2x^4-24\pi x^2+3)e^{-\pi x^2},
\]
有
\[
 \begin{aligned}
 \widehat H(\xi)
 &=\left[
 16\pi^2\left(\xi^4-\frac3\pi\xi^2+\frac3{4\pi^2}\right)
 -24\pi\left(\frac1{2\pi}-\xi^2\right)+3
 \right]q(\xi)\\
 &=(16\pi^2\xi^4-24\pi\xi^2+3)q(\xi)=H(\xi).
 \end{aligned}
\]
这里的自对偶性由两次多项式计算得到，不需要预先使用 Fourier 变换的特征函数理论。

还可以把它的方差完全算出。设
\[
 M_n=\int_{\mathbb R}x^{2n}e^{-2\pi x^2}\,\mathrm dx.
\]
Gaussian 衰减使分部积分的边界项为零，对 \(n\geq1\) 有
\[
 0=\int_{\mathbb R}
 \frac{\mathrm d}{\mathrm dx}
 \left(x^{2n-1}e^{-2\pi x^2}\right)\,\mathrm dx
 =(2n-1)M_{n-1}-4\pi M_n.
\]
结合 \(M_0=1/\sqrt2\)，这给出
\[
 M_n=\frac{(2n-1)!!}{(4\pi)^n\sqrt2},
\]
其中 \(n=0\) 时的空乘积取一。平方多项式为
\[
 (16\pi^2x^4-24\pi x^2+3)^2
 =256\pi^4x^8-768\pi^3x^6
   +672\pi^2x^4-144\pi x^2+9.
\]
逐项代入上述矩公式，得到
\[
 \|H\|_2^2=12\sqrt2,
 \quad
 \int_{\mathbb R}x^2|H(x)|^2\,\mathrm dx
 =\frac{27\sqrt2}{\pi}.
\]
于是归一化函数
\begin{equation}\label{b2:eq:20-nongaussian-selfdual}
 h(x)=\frac{(16\pi^2x^4-24\pi x^2+3)e^{-\pi x^2}}
             {\sqrt{12\sqrt2}}
\end{equation}
满足 \(\widehat h=h\) 和 \(\|h\|_2=1\)。它是偶函数，两中心为零；自对偶性又使两个方差相同，故
\[
 \sigma_x^2=\sigma_\xi^2=\frac9{4\pi},
 \quad
 \sigma_x\sigma_\xi=\frac9{4\pi}>\frac1{4\pi}.
\]
多项式在 \(x=0\) 时为正，在 \(x^2=3/(4\pi)\) 时为负，故 \(h\) 有零点并变号，不属于前述非零 Gaussian 极值族。这个例子把自对偶与极值的区别具体化：Fourier 变换保持函数本身，并不保证这个函数已经具有最小的方差乘积。
